\section{仲裁协议}
\begin{enumerate}
  \item 必考的内容。
  \item \zd{仲裁协议是仲裁制度的基础，是仲裁约束力的来源}。
  \item 形式。
        \begin{enumerate}
          \item \zd{必须书面}，口头无效。可以表现为合同中的一个仲裁条款，也可以表现为单独的一份仲裁协议书。
          \item 包含三个方面内容：
                \begin{enumerate}
                  \item 请求仲裁的意思表示。
                  \item 仲裁事项必须要具有可仲裁性。
                  \item 要有选定的仲裁委员会。\zd{明确具体唯一}。
                        \begin{enumerate}
                          \item 约定不明确仲裁协议无效。
                          \item 约定了两个以上仲裁机构的，可以协议选一个。
                          \item 约定的名称不准确，但能够确定具体仲裁机构的。
                          \item 根据仲裁规则能够确定具体仲裁委的视为约定明确。
                          \item 既可以仲裁又可以起诉这种条款叫\zd{或诉或裁条款}，无效。
                        \end{enumerate}
                \end{enumerate}
        \end{enumerate}
  \item 效力。\zd{一次性，用过即废}。
        \begin{enumerate}
          \item 有效的仲裁协议将排斥法院的司法管辖权（仅仅阻止当事人通过诉讼方式解决纠纷，不排除特别程序）。
                \begin{enumerate}
                  \item 在受理前发现仲裁协议的裁定不予受理。
                  \item 受理后发现的裁定驳回起诉。
                  \item 仲裁协议无效的除外。
                  \item 一方向法院起诉另外一方在法院的首次开庭前提出仲裁协议的，
                        视为放弃仲裁协议接受法院管辖，法院取得应诉管辖。
                \end{enumerate}
          \item 独立性。独立于合同而存在的，不会因为合同的无效而无效，也不会因为合同的解除而解除。
          \item 只对在仲裁协议上签字的双方当事人生效 \yjt 仲裁协议效力相对性。仲裁协议效力的扩张：
                \begin{enumerate}
                  \item 权利义务的继承人。自然人死亡，企业合并分立。
                  \item 权利义务转让的对受让人有约束力，但受让人在受让合同时明确表示反对的 \hz 不知道有单独的仲裁协议的除外。
                \end{enumerate}
          \item 仲裁协议效力的确认。
                \begin{enumerate}
                  \item 两个主体：仲裁委、法院（申请人/被申请人住所地；协议签订地、约定的仲裁机构所在地）。
                  \item 法院确认优先，但仲裁委先于法院接受申请并且已经作出决定的法院就不再受理了。
                  \item 要主张仲裁协议无效只能\zd{在首次开庭前提出}，之后不准再对本案的裁决提出任何的质疑和非难。
                \end{enumerate}
        \end{enumerate}
\end{enumerate}
